\DocumentMetadata{
  pdfversion=2.0,
  pdfstandard=ua-2,
  lang=en-US,
  tagging=on,
  tagging-setup={math/setup=mathml-SE,tabsorder=structure}
}
\documentclass[11pt,letterpaper]{article}
\usepackage[margin=0.68in,includefoot]{geometry}
\usepackage{newcomputermodern}
\usepackage{amsmath,amscd,mathtools}
\usepackage{tikz,adjustbox}
\providecommand{\square}{\mdlgwhtsquare}
\usepackage[hidelinks]{hyperref}
\hypersetup{
  pdftitle={Complex Analysis PhD Exam, September 1988},
  pdfauthor={University of Florida Department of Mathematics},
  pdfsubject={Graduate mathematics examination},
  pdfdisplaydoctitle=true,
  bookmarksopen=true
}
\setcounter{secnumdepth}{0}
\setlength{\parindent}{0pt}
\setlength{\parskip}{0.28em}
\setlength{\emergencystretch}{3em}
\setlength{\leftmargini}{2.15em}
\setlength{\leftmarginii}{2.65em}
\setlength{\leftmarginiii}{3em}
\renewcommand{\labelenumi}{\textbf{\theenumi.}}
\pagestyle{empty}
\raggedbottom
\allowdisplaybreaks
\newenvironment{examproblems}[1]{%
  \begin{enumerate}%
  \setcounter{enumi}{#1}%
  \setlength{\topsep}{0.35em}%
  \setlength{\partopsep}{0pt}%
  \setlength{\itemsep}{0.48em}%
  \setlength{\parsep}{0.18em}%
}{\end{enumerate}}
\newenvironment{examsubparts}{%
  \begin{enumerate}%
  \setlength{\topsep}{0.18em}%
  \setlength{\partopsep}{0pt}%
  \setlength{\itemsep}{0.16em}%
  \setlength{\parsep}{0.1em}%
}{\end{enumerate}}
\newenvironment{examinstructions}{%
  \par\begingroup\itshape
}{%
  \par\endgroup\vspace{0.18em}
}
\makeatletter
\renewcommand\section{\@startsection{section}{1}{0pt}%
  {-1.25ex plus -.25ex minus -.1ex}%
  {0.55ex plus .1ex}%
  {\normalfont\large\bfseries}}
\renewcommand{\@maketitle}{%
  \newpage\null\vskip -1em%
  {\footnotesize\bfseries
    \href{https://www.ufl.edu/}{University of Florida}, \href{https://math.ufl.edu/}{Department of Mathematics}\par}%
  \vskip -0.5em%
  \tagstructbegin{tag=Title}%
    {\LARGE\bfseries\href{https://gma.math.ufl.edu/exams/complex-analysis-phd/}{Complex Analysis PhD Exam}, September 1988\par}%
  \tagstructend%
  \vskip 0.25em%
  \tagmcbegin{artifact}%
    \hrule height 1pt%
  \tagmcend%
  \vskip 1em%
}
\makeatother
\title{Complex Analysis PhD Exam, September 1988}
\author{}
\date{}
\begin{document}
\maketitle
\thispagestyle{empty}
\begin{examinstructions}
Present your work clearly and do not leave any gaps in the proofs. Do all 9 problems.
\end{examinstructions}
\begin{examproblems}{0}
\item
Suppose~\(f\) is one-to-one and entire. Show that \(f(z)=az+b\), \(a\ne 0\).
\item
Let \(J\ne 0\) be a complex number. Prove that if the series \(\sum_{n=0}^{\infty}a_nJ^n\) converges, then the power series \(\sum_{n=0}^{\infty}a_nz^n\) converges for \(|z|<|J|\).
\item
Evaluate 
\[
\int_0^\infty \frac{\cos x}{1+x^2}\,dx.
\]
\item
This problem has two parts.
\begin{examsubparts}
\item[\textnormal{a)}]
Let \(f(z)=\sum_{n=0}^{\infty}z^{2^n}\), \(|z|<1\). Show that the unit circle is a natural boundary for~\(f\).
\item[\textnormal{b)}]
Let \(g(z)=\prod_{n=0}^{\infty}(1+z^{2^n})\). Discuss the behavior of \(g(z)\) on \(|z|=1\).
\end{examsubparts}
\item
Find all entire functions~\(f\) such that \(|f(z)|=1\) whenever \(|z|=1\).
\item
Suppose~\(f\) is analytic in the upper half-plane \(\operatorname{Im}z>0\) and \(|f|\le 1\). How large can \(|f'(i)|\) be? Find the extremal functions.
\item
Suppose \(\{f_n\}\) is a sequence of analytic functions on the region~\(D\), none of the functions \(f_n\) has a zero in~\(D\), and \(\{f_n\}\) converges to~\(f\) uniformly on compact subsets of~\(D\). Prove that either~\(f\) has no zero in~\(D\) or \(f(z)=0\) for all~\(z\) in~\(D\).
\item
This problem has two parts.
\begin{examsubparts}
\item[\textnormal{a)}]
Suppose~\(f\) is analytic in a region~\(D\) containing the closed unit disc \(|z|\le 1\), \(|f(z)|>2\) if \(|z|=1\), and \(f(0)=1\). Must~\(f\) have a zero in the unit disc?
\item[\textnormal{b)}]
Suppose~\(f\) and~\(g\) are entire functions, and \(|f(z)|\le |g(z)|\) for every~\(z\). What conclusion can you draw?
\end{examsubparts}
\item
Suppose~\(f\) is an entire function and \(f(0)=1\). Let 
\[
M(r)=\sup_{0\le t\le 2\pi}|f(re^{it})|
\]
 and \(n(r)\) be the number of zeros of~\(f\) in the disc \(|z|\le r\). Prove that 
\[
n(r)\ln 2\le \ln M(2r).
\]
\end{examproblems}
\end{document}
